%! Properties and Applications of Determinants — Unit 5, Lesson 5.2 — Homework (Answer Key)
\documentclass[10pt]{article}
\usepackage{linalg-article}
\usepackage{linalg-key}   % pulls in -boxes; do NOT also load -boxes
\newcommand{\TallMath}[1]{$\displaystyle #1\rule[-1.4em]{0pt}{3.2em}$}
\newcommand{\vv}{\mathbf{v}}
\newcommand{\ww}{\mathbf{w}}
\newcommand{\avec}{\mathbf{a}}
\newcommand{\bb}{\mathbf{b}}
\newcommand{\xx}{\mathbf{x}}
\newcommand{\T}{\mathsf{T}}
\newcommand{\zero}{\mathbf{0}}

\begin{document}

\pageheader{Unit 5, Lesson 5.2}{Homework --- Answer Key}
\namedateperiod

\vspace{0.12in}

\begin{notesbox}{Practice}
Use the \textbf{rules}, not cofactor expansion: reduce to triangular and multiply the pivots (flip the
sign once per row swap), and remember that determinants \emph{multiply}: $\det(AB)=\det A\,\det B$.

\begin{enumerate}[leftmargin=1.9em, itemsep=11pt]
  \item \textbf{Compute by elimination.}
    \begin{enumerate}[label=(\alph*), leftmargin=1.8em, itemsep=5pt]
      \item $\det\begin{bmatrix} 2 & 1 & 0 \\ 4 & 3 & 1 \\ 0 & 1 & 3 \end{bmatrix}$ \hfill \ans{$R_2-2R_1,\,R_3-R_2\Rightarrow$ pivots $2,1,2$; $\det=4$}
      \item $\det\begin{bmatrix} 3 & 1 & 2 \\ 0 & 2 & 5 \\ 0 & 0 & 4 \end{bmatrix}$ \; {\footnotesize(triangular)} \hfill \ans{$3\cdot2\cdot4=24$}
      \item $\det\begin{bmatrix} 1 & 2 & 3 \\ 2 & 4 & 6 \\ 1 & 0 & 5 \end{bmatrix}$ \hfill \ans{$0$ (row 2 $=2\times$ row 1)}
    \end{enumerate}
  \item \textbf{What does each move do to the determinant?} A matrix $A$ has $\det A=6$. Find the new
        determinant after each single operation.
    \begin{enumerate}[label=(\alph*), leftmargin=1.8em, itemsep=5pt]
      \item Swap two rows. \hfill \ans{$-6$ (swap negates)}
      \item Multiply one row by $5$. \hfill \ans{$30$ (scale by $5$)}
      \item Add $2\times(\text{row }1)$ to row $3$. \hfill \ans{$6$ (shear --- unchanged)}
    \end{enumerate}
  \item \textbf{Product rule \& inverses.} Let $A=\begin{bmatrix} 3 & 1 \\ 0 & 2 \end{bmatrix}$ and
        $B=\begin{bmatrix} 2 & 0 \\ 1 & 4 \end{bmatrix}$. Compute $\det A$ and $\det B$, then find
        $\det(AB)$ (two ways: the product rule, and by multiplying $AB$ out). What is $\det A^{-1}$?
        \ansline{$\det A=6$, $\det B=8$. Product rule: $\det(AB)=6\cdot8=48$. Directly, $AB=\begin{bsmallmatrix} 7 & 4\\ 2 & 8\end{bsmallmatrix}$, $\det=56-8=48$ --- matches. $\det A^{-1}=1/\det A=1/6$.}
  \item \textbf{Explain.} A $3\times3$ matrix reduces to a triangular matrix with a $0$ on the diagonal.
        What is its determinant? Is the matrix invertible? Explain in one sentence.
        \ansline{$\det=0$ (the product of the diagonal includes a $0$), so the matrix is \emph{not} invertible --- $\det=0$ means the box is flat, so there is no inverse.}
\end{enumerate}
\end{notesbox}

\begin{extensionbox}
\textbf{Cramer's rule: solving a system with determinants.} For $A\xx=\bb$, each unknown is a ratio of
determinants --- replace a column of $A$ with $\bb$, take the determinant, and divide by $\det A$.
\begin{enumerate}[label=(\alph*), leftmargin=1.8em, itemsep=5pt]
  \item Solve $\begin{cases} \ x + 2y = 4 \\ 3x + \ y = 7 \end{cases}$ with
        $A=\begin{bsmallmatrix} 1 & 2 \\ 3 & 1 \end{bsmallmatrix}$, $\bb=\begin{bsmallmatrix} 4 \\ 7 \end{bsmallmatrix}$.
        Compute $\det A$, then $x=\det A_x/\det A$ and $y=\det A_y/\det A$.
  \item \textbf{Transpose rule.} For $A=\begin{bsmallmatrix} 4 & 1 \\ 2 & 3 \end{bsmallmatrix}$, check that $\det A^{\T}=\det A$.
\end{enumerate}
\ansline{(a) $\det A=(1)(1)-(2)(3)=-5$. $A_x=\begin{bsmallmatrix} 4 & 2\\ 7 & 1\end{bsmallmatrix}$, $\det A_x=4-14=-10$, so $x=-10/-5=2$. $A_y=\begin{bsmallmatrix} 1 & 4\\ 3 & 7\end{bsmallmatrix}$, $\det A_y=7-12=-5$, so $y=-5/-5=1$. Check: $2+2(1)=4$ \checkmark, $3(2)+1=7$ \checkmark.}
\ansline{(b) $\det A=(4)(3)-(1)(2)=10$; $A^{\T}=\begin{bsmallmatrix} 4 & 2\\ 1 & 3\end{bsmallmatrix}$, $\det A^{\T}=12-2=10$. Equal.}
\end{extensionbox}

\begin{spiralbox}
\textbf{Coming next --- Lesson 5.3: Linear Transformations.} A matrix doesn't just hold numbers --- it
\emph{moves space}. Next we watch what $A$ does to shapes, and see the determinant reappear as the
\textbf{factor every area and volume is multiplied by}, with its \emph{sign} reporting whether
orientation was kept or flipped.
\end{spiralbox}

\begin{teachernote}
Item~1 practices the fast method: full elimination (a), the triangular shortcut (b), and a dependent
row $\Rightarrow0$ (c). Item~2 isolates the three row rules --- \emph{swap negates}, \emph{scale
multiplies}, \emph{shear does nothing} --- the single most common point of confusion. Item~3 is the
product rule two ways plus $\det A^{-1}=1/\det A$; item~4 the $\det=0\Leftrightarrow$ singular
justification. The extension is the application layer: Cramer's rule (watch signs when a determinant is
negative) and the transpose rule. Most common slips: treating a shear like a scale, dropping the sign
on a swap, and the false ``$\det(A+B)=\det A+\det B$.''
\end{teachernote}

\end{document}
